\chapter*{Conclusion générale}
\addcontentsline{toc}{chapter}{Conclusion générale}

Ce projet de fin d'études s'est inscrit dans un contexte où la transformation numérique pousse les organisations à consolider des données RH et projet de plus en plus volumineuses, tout en exigeant des décisions plus rapides, plus traçables et mieux partagées entre les acteurs. Au sein de l'écosystème \textbf{Augmented Talents}, le travail réalisé a porté sur \textbf{Strategic Copilot} : une plateforme SaaS d'aide à la décision pour la gestion stratégique des talents et du portefeuille projet. L'ambition était de rapprocher, dans une interface unique, la visualisation des indicateurs, la détection des risques, l'analyse de viabilité des projets, les recommandations assistées par intelligence artificielle et la gestion structurée des demandes RH, pour les profils manager, responsable RH et talent. Encadré pédagogiquement par Madame Soundouss Abroun à l'École Marocaine des Sciences de l'Ingénieur de Tanger, et professionnellement par Monsieur Youssef Moufak ainsi que par Messieurs Hamza Sadouk, Hicham Taoufik et Mohyi Eddine Saadani, le stage a visé à livrer une solution opérationnelle, fondée sur le dépôt \texttt{frontend-ui} et sur le document de cadrage \emph{PFE Agentic AI Strategic Talents}.

Le rapport a retracé cette démarche en six chapitres complémentaires. Après le contexte et l'analyse des besoins, l'étude bibliographique a posé les fondements RH, projet, aide à la décision et technologies web. Le cahier des charges et la planification ont formalisé les spécifications et le calendrier du PFE. La conception a décrit l'architecture multi-agents, les workflows n8n et les API. La réalisation a détaillé l'implémentation React/Vite et les interfaces livrées. Enfin, le chapitre consacré aux tests a explicité la stratégie de validation et les limites de la preuve automatisée. Cette conclusion générale en tire le bilan global.

Sur le plan des \textbf{résultats}, les objectifs fixés en début de projet ont été atteints dans une large mesure sur l'environnement de préproduction connecté à \texttt{n8nprod.aphelionxinnovations.com}. L'application \texttt{copilot-ui} propose trois workspaces cohérents (\texttt{/workspace/manager}, \texttt{/workspace/rh}, \texttt{/workspace/talent}), protégés par authentification JWT et contrôle d'accès par rôle. Le manager dispose d'un tableau de bord, d'un portefeuille projets, du modal \emph{Mission Control} (pilotage à 360 degrés : overview, équipe, tâches, risques, simulation, décisions), de la gestion des alertes, des demandes RH, du Helper conversationnel, du simulateur what-if et du journal des décisions enrichi de graphiques Recharts. Le responsable RH accède aux agrégats métier, au référentiel talents et compétences, aux plans de formation, à la mobilité, aux alertes organisationnelles et à la gestion des comptes. Le talent consulte ses missions, tâches, charge, compétences et formations via des endpoints dédiés \texttt{/api/workspace/talent/*}.

L'architecture retenue respecte le principe \textbf{API-first} et la règle, inscrite dans \texttt{crud-domain.ts}, selon laquelle le front-end n'effectue aucun recalcul de scores de viabilité, de risque ou de recommandations : ces traitements sont délégués aux workflows n8n, qui orchestrent les agents documentés (analyse projet, matching talents, risques et signaux faibles, orchestrateur de viabilité, Helper Copilot) et persistent les résultats dans PostgreSQL, hébergé via Supabase. Les scores et décisions \texttt{Continue}, \texttt{Adjust} et \texttt{Stop} sont produits côté serveur ; l'utilisateur conserve la responsabilité de la validation, notamment via le Strategist, les demandes RH et l'enregistrement des décisions Copilot. Les campagnes de validation menées pendant le stage --- recettes manuelles des parcours critiques, contrôles \texttt{npm run typecheck} et \texttt{npm run build:check}, test unitaire des traces d'agents --- confirment la cohérence fonctionnelle de cette chaîne sur les données et workflows disponibles en préproduction.

Ce travail a permis de consolider des \textbf{compétences techniques} variées. Sur le plan front-end, la maîtrise de React 19, TypeScript, Vite, Tailwind CSS 4, React Router 7 et TanStack React Query s'est traduite par une SPA multi-rôles maintenable, avec internationalisation (français, anglais, arabe) et accessibilité via React Aria. L'intégration de dizaines de webhooks versionnés (\texttt{wmp-*-v1}, \texttt{wmt-*-v1}, routes \texttt{/webhook/api/*}) a exigé une rigueur sur les contrats JSON, les parseurs défensifs et la configuration par variables \texttt{VITE\_*}. La mise en place du proxy de développement et du relais Vercel a renforcé la compréhension des enjeux CORS, de sécurité JWT et de déploiement SaaS.

Sur le plan \textbf{méthodologique}, ce stage m'a permis de mener un projet complet, de l'analyse des besoins à la validation : cadrage agentique (Analyst, Matchmaker, Watchdog, Orchestrator), conception de l'architecture, développement frontend et backend, modélisation de la base de données, intégration de l'IA et recettes sur l'environnement de préproduction. La rédaction du présent rapport en LaTeX modulaire a structuré la communication écrite de ce travail. Les échanges réguliers avec les encadrants ont contribué à aligner les choix techniques sur des objectifs métier explicites.

Plusieurs \textbf{limites} demeurent toutefois assumées. La couverture de tests automatisés reste faible dans le dépôt (essentiellement \texttt{agent-trace.test.ts}), ce qui rend les régressions dépendantes de recettes manuelles. L'environnement n8n et les données de préproduction partagés limitent la reproductibilité stricte des scores lorsque l'IA générative est activée (\texttt{use\_ai: true}). La latence de certains workflows impose des timeouts élevés et peut dégrader l'expérience utilisateur. La multiplicité des conventions d'URL (chemins versionnés, webhooks UUID) complexifie la configuration et la maintenance. Enfin, la qualité des recommandations et des alertes reste conditionnée par l'exhaustivité et la fraîcheur des données en base, facteur hors du seul front-end mais déterminant pour l'adoption terrain.

Les \textbf{perspectives} d'évolution s'articulent naturellement autour de ces limites. À court terme, l'extension des tests automatisés (Playwright, tests unitaires sur les parseurs API) et l'archivage d'une spécification OpenAPI ou d'une collection Postman sécuriseraient les livraisons. À moyen terme, un environnement de démonstration avec jeu de données figé, des indicateurs de progression plus fins pendant les analyses longues et une harmonisation des chemins n8n simplifieraient l'exploitation. Sur le plan produit Augmented Talents, l'enrichissement des scénarios what-if, le renforcement de la gouvernance des décisions IA et l'élargissement des rapports PDF automatisés pourraient accroître la valeur perçue par les directions RH et opérationnelles. L'intégration continue entre workflows n8n et interface React demeure le fil conducteur : faire évoluer les agents sans remettre en cause l'expérience utilisateur, grâce à des contrats stables et à une traçabilité exemplaire.

En définitive, Strategic Copilot démontre la faisabilité d'un copilote stratégique agentique au service du pilotage conjoint des talents et des projets. Le livrable \texttt{copilot-ui} matérialise une architecture moderne, découplée et orientée aide à la décision responsable, où l'humain reste au centre des arbitrages. Pour l'auteur de ce mémoire, Ouissame Mahtour, ce projet de fin d'études constitue une synthèse exigeante entre ingénierie logicielle, gestion des ressources humaines et intelligence artificielle appliquée ; il ouvre sur une contribution durable à la plateforme Augmented Talents et sur une pratique professionnelle renforcée par la rigueur, la curiosité technique et le souci de la valeur métier.
\newpage